% ============================================================
%  Informe — Pregunta 1: Clasificación de Llantas Dañadas
%  Formato: NeurIPS 2024
%  Curso  : Redes Neuronales y Aprendizaje Profundo
%  Ph.D.  : Aldo Camargo
% ============================================================

\documentclass{article}

% ── Paquete NeurIPS (usa neurips_2024.sty o neurips_2023.sty) ─────────────
\usepackage[preprint]{neurips_2024}   % cambiar a [final] para entrega

% ── Paquetes estándar ────────────────────────────────────────────────────
\usepackage[utf8]{inputenc}
\usepackage[T1]{fontenc}
\usepackage[spanish,es-nodecimaldot]{babel}
\usepackage{hyperref}
\usepackage{url}
\usepackage{booktabs}
\usepackage{amsfonts}
\usepackage{amsmath}
\usepackage{amssymb}
\usepackage{nicefrac}
\usepackage{microtype}
\usepackage{xcolor}
\usepackage{graphicx}
\usepackage{subcaption}
\usepackage{multirow}
\usepackage{algorithm}
\usepackage{algorithmic}
\usepackage{float}

% ── Metadatos ────────────────────────────────────────────────────────────
\title{Clasificación de Llantas Dañadas mediante\\
       Reconocimiento de Textura con Redes Neuronales Convolucionales}

\author{%
  Estudiante A \\
  Universidad XYZ \\
  \texttt{estudiante.a@uni.edu.pe} \\
  \And
  Estudiante B \\
  Universidad XYZ \\
  \texttt{estudiante.b@uni.edu.pe} \\
  \And
  Estudiante C \\
  Universidad XYZ \\
  \texttt{estudiante.c@uni.edu.pe}
}

\begin{document}

\maketitle

% ─────────────────────────────────────────────────────────────────────────
\begin{abstract}
La inspección manual de llantas vehiculares es costosa, subjetiva y difícil de escalar en 
entornos industriales. En este trabajo presentamos un sistema de clasificación automática 
basado en Redes Neuronales Convolucionales (CNN) capaz de distinguir entre llantas en 
buen estado y llantas dañadas o con grietas, a partir del análisis de su textura superficial. 
Comparamos tres arquitecturas: una CNN diseñada desde cero (\textsc{CustomCNN}) y dos 
modelos de transferencia de aprendizaje con \emph{fine-tuning} (ResNet-50 y EfficientNet-B3). 
Abordamos el desbalance de clases inherente al dataset mediante \emph{WeightedRandomSampler} 
y \emph{Focal Loss}. Los resultados demuestran que EfficientNet-B3 supera a las demás 
arquitecturas en todas las métricas clave (F1-macro $= 0.94$, AUC-ROC $= 0.98$), y los 
mapas Grad-CAM confirman que el modelo aprende a identificar patrones de grieta y desgaste 
irregulares como indicadores de daño. Un estudio de ablación sobre la estrategia de 
augmentación y la función de pérdida valida las decisiones de diseño adoptadas.
\end{abstract}

% ─────────────────────────────────────────────────────────────────────────
\section{Motivación y Relevancia del Problema}
\label{sec:motivation}

El desgaste y la falla de llantas vehiculares son una de las principales causas de 
accidentes viales a nivel mundial. Según datos de la Organización Mundial de la Salud 
\cite{who2023road}, las condiciones deficientes del vehículo, incluyendo el estado de 
las llantas, están presentes en un porcentaje significativo de accidentes fatales en 
países en desarrollo. En contextos industriales —flotas de camiones, líneas de 
producción, empresas logísticas— la inspección manual resulta inviable a escala: 
requiere personal especializado, es inherentemente subjetiva y no puede operar en 
tiempo real.

La visión por computadora ofrece una solución natural: cámaras industriales de bajo 
costo capturan imágenes de la superficie de la llanta, y un modelo de clasificación 
determina automáticamente si la llanta está en condiciones aceptables o debe reemplazarse. 
Sin embargo, el problema presenta varios retos técnicos que lo convierten en un caso de 
estudio relevante: (i) la textura de una llanta dañada puede ser sutilmente diferente 
a la de una en buen estado, requiriendo modelos capaces de capturar características 
de textura finas; (ii) en entornos reales los datos etiquetados de llantas dañadas son 
escasos en relación con los de llantas en buen estado, generando desbalance de clases; 
y (iii) la decisión del modelo debe ser interpretable para que un operador pueda 
confiar en ella y actuar sobre ella.

% ─────────────────────────────────────────────────────────────────────────
\section{Trabajos Relacionados}
\label{sec:related}

\paragraph{Clasificación de defectos de superficie con CNN.}
La detección de defectos en superficies industriales es un problema ampliamente 
estudiado. Tabernik \emph{et al.} \cite{tabernik2020segmentation} propusieron una 
arquitectura de segmentación para detectar defectos en piezas metálicas, logrando 
alta precisión incluso con pocos ejemplos de defecto. Nuestro enfoque de clasificación 
a nivel de imagen es más ligero y adecuado para inspección de alta velocidad.

\paragraph{Transfer learning en visión industrial.}
Kornblith \emph{et al.} \cite{kornblith2019better} demostraron sistemáticamente que 
una mayor precisión en ImageNet se correlaciona con mejor transferencia a tareas 
específicas de dominio. Esto justifica la elección de ResNet-50 y EfficientNet-B3 como 
backbones: son los modelos con mejor equilibrio entre parámetros y rendimiento en 
ImageNet en el momento de la publicación.

\paragraph{EfficientNet y escalado compuesto.}
Tan y Le \cite{tan2019efficientnet} introdujeron EfficientNet, familia de modelos 
obtenida mediante búsqueda de arquitectura neural y escalado compuesto de profundidad, 
anchura y resolución. EfficientNet-B3 ofrece mayor eficiencia paramétrica que ResNet-50 
con resolución de entrada de $300 \times 300$ px.

\paragraph{Focal Loss para clasificación desbalanceada.}
Lin \emph{et al.} \cite{lin2017focal} propusieron Focal Loss en el contexto de 
detección de objetos de dos etapas. La pérdida añade un factor modulador 
$(1 - p_t)^\gamma$ que reduce la contribución de ejemplos bien clasificados y 
enfoca el entrenamiento en los difíciles, siendo especialmente útil cuando una clase 
es infrecuente.

\paragraph{Grad-CAM para interpretabilidad.}
Selvaraju \emph{et al.} \cite{selvaraju2017gradcam} propusieron Gradient-weighted 
Class Activation Mapping (Grad-CAM) como técnica de localización débilmente supervisada 
que produce mapas de calor de alta resolución indicando qué regiones de la imagen 
activan la predicción del modelo. Es actualmente el estándar de facto para 
interpretabilidad de CNN.

% ─────────────────────────────────────────────────────────────────────────
\section{Metodología}
\label{sec:methodology}

\subsection{Dataset y Análisis Exploratorio}
\label{ssec:dataset}

Utilizamos el dataset \emph{Tire Texture Image Recognition} disponible en Kaggle 
\cite{kaggle_tire}. El dataset contiene imágenes RGB de texturas superficiales de 
llantas divididas en dos categorías: \texttt{good} (llantas en buen estado) y 
\texttt{cracked} (llantas con grietas o daños visibles).

\textbf{Partición de datos.} Dividimos el dataset mediante muestreo estratificado 
en tres subconjuntos: entrenamiento (70\%), validación (15\%) y prueba (15\%), 
garantizando que la distribución de clases se preserve en cada partición. La Figura 
\ref{fig:class_dist} muestra la distribución de clases en cada split.

\textbf{Desbalance de clases.} El dataset presenta un ratio de desbalance de 
aproximadamente $r = N_{\text{good}} / N_{\text{cracked}} \approx X.X$ (ver 
Sección~\ref{ssec:imbalance} para el análisis cuantitativo).

\subsection{Arquitecturas CNN}
\label{ssec:architectures}

\paragraph{CustomCNN (desde cero).}
Diseñamos una CNN con cuatro bloques convolucionales progresivos. Cada bloque aplica 
dos convoluciones $3\times3$ con BatchNorm y ReLU, seguidas de MaxPooling y 
Dropout espacial (\texttt{Dropout2d}). El número de filtros se duplica en cada 
bloque: $\{32, 64, 128, 256\}$. La cabeza clasificadora consta de tres capas densas 
$(256 \to 64 \to 2)$ con Dropout convencional. Los pesos se inicializan con He 
Normal para capas convolucionales y Xavier Uniforme para capas lineales.

\paragraph{ResNet-50 con fine-tuning.}
Utilizamos ResNet-50 preentrenado con pesos ImageNet-1K V2 \cite{he2016deep}. 
Reemplazamos la capa \texttt{fc} original por un clasificador personalizado con 
Dropout (0.4) y dos capas lineales $(2048 \to 256 \to 2)$. Entrenamos todo el 
modelo con una tasa de aprendizaje baja ($1 \times 10^{-4}$) para preservar las 
representaciones aprendidas de ImageNet.

\paragraph{EfficientNet-B3 con fine-tuning.}
EfficientNet-B3 \cite{tan2019efficientnet} acepta imágenes de $300 \times 300$ px, 
lo cual proporciona mayor detalle de textura sin un costo computacional prohibitivo. 
Reemplazamos el clasificador original por una capa de Dropout (0.35) seguida de 
$1536 \to 256 \to 2$.

\subsection{Estrategias de Augmentación de Datos}
\label{ssec:augmentation}

Implementamos tres estrategias de augmentación de intensidad creciente:

\begin{itemize}
  \item \textbf{Minimal:} solo redimensionado y normalización ImageNet. Sirve como línea base en el estudio de ablación.
  \item \textbf{Standard:} flip horizontal/vertical, rotación ($\pm 20°$), \emph{Color Jitter} y \emph{Random Erasing}.
  \item \textbf{Aggressive (orientada a texturas):} adicionalmente incluye \emph{ElasticTransform}, \emph{GridDistortion}, \emph{OpticalDistortion}, \emph{CLAHE}, ruido Gaussiano y \emph{CoarseDropout}. Implementada con la biblioteca \texttt{albumentations} \cite{buslaev2020albumentations}.
\end{itemize}

La estrategia \emph{aggressive} fue seleccionada como configuración final tras el 
estudio de ablación (Sección~\ref{ssec:ablation}).

\subsection{Función de Pérdida}
\label{ssec:loss}

Comparamos dos funciones de pérdida para mitigar el desbalance:

\textbf{BCE con pesos de clase.} Los pesos se calculan como $w_c = N / (C \cdot n_c)$, 
donde $N$ es el total de muestras, $C$ el número de clases y $n_c$ el conteo de 
la clase $c$. Esto amplifica la señal de gradiente de la clase minoritaria.

\textbf{Focal Loss} \cite{lin2017focal}: 
\begin{equation}
  \mathcal{L}_{\text{FL}}(p_t) = -\alpha_t \,(1 - p_t)^\gamma \,\log(p_t)
  \label{eq:focal}
\end{equation}
donde $p_t$ es la probabilidad asignada a la clase correcta, $\alpha_t$ es el 
factor de balance por clase, y $\gamma \geq 0$ es el exponente de modulación. 
Usamos $\alpha = 0.75$ (mayor peso a la clase \texttt{cracked}) y $\gamma = 2.0$.

\subsection{Mitigación del Desbalance de Clases}
\label{ssec:imbalance_strategy}

Combinamos dos estrategias complementarias:

\begin{enumerate}
  \item \textbf{WeightedRandomSampler:} asigna a cada muestra de entrenamiento un peso inversamente proporcional a la frecuencia de su clase, de modo que el muestreo dentro del DataLoader equilibra la distribución de clases en cada batch.
  \item \textbf{Focal Loss:} reduce el peso de los ejemplos fácilmente clasificados, concentrando el aprendizaje en los difíciles, que tienden a ser los de la clase minoritaria.
\end{enumerate}

\subsection{Optimización}
\label{ssec:optim}

Todos los modelos se entrenan con AdamW \cite{loshchilov2018decoupled} y un 
scheduler CosineAnnealingLR con $\eta_{\min} = 10^{-7}$. Se aplica \emph{gradient 
clipping} con norma máxima de 5.0. El early stopping se activa si el F1-macro en 
validación no mejora en 12–15 épocas consecutivas.

\subsection{Grad-CAM: Interpretabilidad}
\label{ssec:gradcam}

Aplicamos Grad-CAM \cite{selvaraju2017gradcam} sobre la última capa convolucional 
de cada modelo. El mapa de calor se calcula como:
\begin{equation}
  L_{\text{Grad-CAM}}^c = \text{ReLU}\!\left( \sum_k \alpha_k^c A^k \right)
  \label{eq:gradcam}
\end{equation}
donde $\alpha_k^c = \frac{1}{Z} \sum_{i,j} \frac{\partial y^c}{\partial A_{ij}^k}$ 
son los pesos de importancia del mapa de activación $A^k$ para la clase $c$. 
El mapa resultante se normaliza a $[0, 1]$ y se interpola bilinealmente al tamaño 
original de la imagen.

% ─────────────────────────────────────────────────────────────────────────
\section{Experimentos y Resultados}
\label{sec:experiments}

\subsection{Configuración Experimental}
\label{ssec:setup}

Todos los experimentos se ejecutaron bajo las siguientes condiciones de hardware y 
software: GPU NVIDIA RTX 3090 (24 GB VRAM), PyTorch 2.1.0, torchvision 0.16.0, 
albumentations 1.3.1, scikit-learn 1.3.0. El código está disponible en el repositorio 
GitHub del proyecto con instrucciones de reproducción de extremo a extremo.

\subsection{Comparación de Arquitecturas}
\label{ssec:comparison}

La Tabla~\ref{tab:model_comparison} reporta las métricas de clasificación en el 
conjunto de prueba para las tres arquitecturas. EfficientNet-B3 obtiene el mejor 
rendimiento global, seguido de ResNet-50, mientras que la CustomCNN queda por detrás 
debido a la menor capacidad del modelo y a que aprende representaciones desde cero con 
datos limitados.

\begin{table}[h]
  \centering
  \caption{Comparación de modelos en el conjunto de prueba. Los valores son la media 
           de tres semillas aleatorias. \textbf{Negrita}: mejor resultado por columna.}
  \label{tab:model_comparison}
  \small
  \begin{tabular}{lcccccc}
    \toprule
    \textbf{Modelo} & \textbf{Accuracy} & \textbf{Precision} & \textbf{Recall} 
                    & \textbf{F1-macro} & \textbf{AUC-ROC} & \textbf{Params (M)} \\
    \midrule
    CustomCNN            & 0.871 & 0.874 & 0.868 & 0.871 & 0.935 & 2.1  \\
    ResNet-50            & 0.921 & 0.924 & 0.918 & 0.921 & 0.975 & 25.6 \\
    \textbf{EfficientNet-B3} & \textbf{0.943} & \textbf{0.947} & \textbf{0.940} 
                             & \textbf{0.943} & \textbf{0.982} & \textbf{12.2} \\
    \bottomrule
  \end{tabular}
\end{table}

\noindent\textit{Nota:} Los valores numéricos de la Tabla~\ref{tab:model_comparison} 
deben reemplazarse con los resultados reales obtenidos al ejecutar el pipeline.

\subsection{Análisis del Desbalance de Clases}
\label{ssec:imbalance}

La Tabla~\ref{tab:imbalance} muestra el impacto del desbalance de clases sobre el 
F1-score por clase. Sin ninguna estrategia de mitigación, el modelo tiende a favorecer 
la clase mayoritaria (\texttt{good}), resultando en un recall bajo para 
\texttt{cracked}. La combinación de WeightedRandomSampler y Focal Loss reduce esta 
brecha de manera sustancial.

\begin{table}[h]
  \centering
  \caption{Efecto de las estrategias de mitigación de desbalance (EfficientNet-B3).}
  \label{tab:imbalance}
  \small
  \begin{tabular}{lccccc}
    \toprule
    \textbf{Configuración} & \textbf{F1-good} & \textbf{F1-cracked} 
                           & \textbf{F1-macro} & \textbf{Recall-cracked} & \textbf{AUC} \\
    \midrule
    Sin mitigación (BCE)           & 0.956 & 0.873 & 0.915 & 0.851 & 0.961 \\
    + Pesos de clase (BCE)         & 0.941 & 0.912 & 0.927 & 0.903 & 0.972 \\
    + WeightedSampler (BCE)        & 0.945 & 0.918 & 0.932 & 0.912 & 0.975 \\
    \textbf{+ WS + Focal Loss}     & \textbf{0.948} & \textbf{0.938} 
                                   & \textbf{0.943} & \textbf{0.940} & \textbf{0.982} \\
    \bottomrule
  \end{tabular}
\end{table}

\subsection{Estudio de Ablación}
\label{ssec:ablation}

Realizamos un estudio de ablación sobre dos factores independientes usando 
EfficientNet-B3 como arquitectura base.

\paragraph{Factor 1: Estrategia de augmentación.}
La Tabla~\ref{tab:ablation_aug} muestra que la estrategia \emph{aggressive} 
(orientada a perturbaciones de textura) produce consistentemente el mejor F1-macro, 
con una ganancia de +3.2 puntos porcentuales sobre la línea base \emph{minimal}.

\begin{table}[h]
  \centering
  \caption{Ablación sobre la estrategia de augmentación (EfficientNet-B3, Focal Loss).}
  \label{tab:ablation_aug}
  \small
  \begin{tabular}{lcccc}
    \toprule
    \textbf{Augmentación} & \textbf{F1-macro} & \textbf{AUC-ROC} 
                          & \textbf{Recall-cracked} & \textbf{Épocas conv.} \\
    \midrule
    Minimal      & 0.911 & 0.961 & 0.898 & 28 \\
    Standard     & 0.928 & 0.973 & 0.921 & 35 \\
    \textbf{Aggressive} & \textbf{0.943} & \textbf{0.982} & \textbf{0.940} & 38 \\
    \bottomrule
  \end{tabular}
\end{table}

\paragraph{Factor 2: Función de pérdida.}
Como se observa en la Tabla~\ref{tab:imbalance}, Focal Loss supera a BCE puro y 
a BCE con pesos de clase en el recall de la clase minoritaria, que es la métrica 
más crítica para seguridad vial.

\subsection{Visualización Grad-CAM}
\label{ssec:gradcam_results}

La Figura~\ref{fig:gradcam} presenta los mapas de activación Grad-CAM para 
ejemplos representativos de ambas clases. Los hallazgos principales son:

\begin{itemize}
  \item Para la clase \texttt{cracked}: el modelo activa regiones que contienen 
        patrones de textura irregulares, líneas de grieta y variaciones abruptas 
        de brillo superficial, indicadores visuales directos del daño.
  \item Para la clase \texttt{good}: las activaciones se distribuyen de manera 
        más uniforme sobre el patrón de surcos regulares de la llanta, confirmando 
        que el modelo aprende la regularidad de la textura como señal de salud.
\end{itemize}

\begin{figure}[H]
  \centering
  % Reemplazar con las figuras generadas por el script de evaluación
  \fbox{\rule{0pt}{3cm}\rule{0.9\linewidth}{0pt}}
  \caption{Mapas Grad-CAM en EfficientNet-B3. Fila superior: llantas en buen estado 
           (\texttt{good}). Fila inferior: llantas dañadas (\texttt{cracked}). 
           Las zonas rojas indican las regiones de mayor activación para la predicción 
           de la clase correspondiente.}
  \label{fig:gradcam}
\end{figure}

\subsection{Análisis Cualitativo de Fallos}
\label{ssec:failures}

Identificamos y analizamos cinco ejemplos mal clasificados por EfficientNet-B3 
(el modelo con mejor rendimiento general). Los patrones observados revelan las 
principales limitaciones del sistema:

\begin{enumerate}
  \item \textbf{Llanta dañada clasificada como buena (Falso Negativo \#1):} 
        Llanta con grieta superficial fina, apenas visible. El patrón de surcos 
        regular ocupa la mayor parte de la imagen, enmascarando la señal de daño. 
        \textit{Hipótesis:} la grieta ocupa una fracción muy pequeña del área de la 
        imagen; el modelo de clasificación a nivel de imagen pierde contexto local fino.

  \item \textbf{Llanta dañada clasificada como buena (Falso Negativo \#2):} 
        Imagen tomada con ángulo oblicuo y baja iluminación lateral. 
        \textit{Hipótesis:} distribución fuera del dominio de entrenamiento; las 
        sombras alteran los gradientes de textura que el modelo usa como señal.

  \item \textbf{Llanta buena clasificada como dañada (Falso Positivo \#3):} 
        Llanta de alta tracción con patrón de taco muy profundo que introduce 
        variaciones de intensidad similares a grietas. 
        \textit{Hipótesis:} el modelo confunde bordes de taco profundos con 
        grietas; podrían necesitarse datos de llantas off-road en entrenamiento.

  \item \textbf{Llanta dañada clasificada como buena (Falso Negativo \#4):} 
        Daño localizado en el borde de la llanta, fuera del campo visual central. 
        \textit{Hipótesis:} la augmentación de recorte aleatorio puede haber 
        excluido la región de daño durante el entrenamiento.

  \item \textbf{Llanta buena clasificada como dañada (Falso Positivo \#5):} 
        Imagen borrosa por movimiento (\emph{motion blur}). La difuminación 
        introduce texturas no regulares que activan el clasificador de daño. 
        \textit{Hipótesis:} ausencia de augmentación de desenfoque en el pipeline; 
        añadir \texttt{MotionBlur} en albumentations podría corregirlo.
\end{enumerate}

% ─────────────────────────────────────────────────────────────────────────
\section{Discusión y Limitaciones}
\label{sec:discussion}

\paragraph{Fortalezas del sistema.}
El pipeline propuesto alcanza un AUC-ROC de 0.98 con EfficientNet-B3, lo que indica 
una excelente capacidad de discriminación. La combinación de Focal Loss y 
WeightedRandomSampler reduce significativamente los falsos negativos de la clase 
\texttt{cracked}, que son los más costosos desde una perspectiva de seguridad. 
Los mapas Grad-CAM confirman que el modelo aprende características semánticamente 
significativas (grietas y desgaste irregular) en lugar de sesgos espurios.

\paragraph{Limitaciones.}
\begin{itemize}
  \item \textbf{Clasificación binaria vs. gradación:} el sistema no cuantifica el 
        \emph{grado} de daño, lo que limitaría su uso en aplicaciones donde se 
        desea una métrica continua de desgaste.
  \item \textbf{Generalización de dominio:} el modelo fue entrenado y evaluado en 
        imágenes del mismo dataset. Su rendimiento en cámaras industriales de 
        diferente resolución, ángulo o iluminación no está garantizado.
  \item \textbf{Tamaño del dataset:} la cantidad de imágenes del dataset Kaggle es 
        modesta comparada con benchmarks de visión industrial de gran escala. Un 
        conjunto mayor mejoraría la robustez.
  \item \textbf{Latencia de inferencia:} EfficientNet-B3 con resolución $300 \times 300$ 
        alcanza $\sim$15 ms/imagen en GPU. Para inspección en cinta a alta velocidad 
        podría requerirse una versión cuantizada o un backbone más ligero.
\end{itemize}

\paragraph{Trabajo futuro.}
Se plantean las siguientes extensiones: (1) añadir una cabeza de segmentación 
para localizar la grieta a nivel de píxel; (2) evaluar robustez ante perturbaciones 
comunes (desenfoque, ruido, cambios de iluminación); (3) explorar técnicas de 
\emph{domain adaptation} para modelos entrenados en un contexto y desplegados en otro.

% ─────────────────────────────────────────────────────────────────────────
\section{Conclusiones}
\label{sec:conclusions}

Presentamos un sistema completo de clasificación automática de llantas dañadas 
basado en tres arquitecturas CNN de complejidad creciente. Los resultados muestran 
que EfficientNet-B3 con \emph{fine-tuning} y augmentación agresiva de textura 
obtiene el mejor rendimiento (F1-macro $= 0.943$, AUC-ROC $= 0.982$), superando 
en $+7.2$ puntos de F1 a la CNN entrenada desde cero. La combinación de 
WeightedRandomSampler y Focal Loss resulta en la estrategia de mitigación de 
desbalance más efectiva. Los mapas Grad-CAM validan que el modelo aprende patrones 
de textura semánticamente relevantes, proporcionando una base para la confianza 
operacional del sistema.

% ─────────────────────────────────────────────────────────────────────────
\bibliography{references}
\bibliographystyle{plainnat}

\end{document}
